% EVAL FIXTURE — see b1.tex. Adequate at stage 6, so the review has no excuse to
% drop below stage 5.
\documentclass[11pt,a4paper]{article}
\usepackage[b2]{ercformatting}
\usepackage{ercreview}

\begin{document}
\maketitle
\startsectiona

\section{Project design}
Eight folds, chosen so that active-site electrostatics and fold ancestry vary
independently -- the contrast O2 rests on. WP1 establishes the ceiling, WP2 attributes
it, WP3 attempts to break it.

\startsectionb

\subsection{WP1: Is the ceiling general?}
\emph{Objective.} A paired-substrate scan for each of the eight folds, with the ceiling
value and its confidence interval.
\emph{Rationale.} O1 is a claim about reproducibility, which two incidental
observations cannot support.

\paragraph{Task 1.1: Paired-substrate scanning.}
To measure both activities on the same variants we run the paired assay across all
eight libraries. Running the pairs in one plate removes the between-assay drift that
made the original saturation look incidental. This yields eight ceiling estimates, the
input to the attribution in WP2.

\subsection{WP2: What sets it?}
\emph{Objective.} An attribution of the ceiling to electrostatics or to ancestry.
\emph{Rationale.} The two candidates co-vary in all published data, so O2 needs the
decorrelated fold set from WP1.

\paragraph{Task 2.3: Moonshot -- raising the ceiling.}
Whether a designed active site can exceed the ceiling is unknown, and a negative result
is informative: it would establish the bound as physical rather than historical, which
is O3's claim either way. WP2's objectives are met by the attribution alone, so the
work package is complete without this task.

\subsection{Risk assessment}
Library quality is the dominant risk: a poorly covered library yields a ceiling
estimate too wide to compare. Coverage is verified before assay, and a failed library
is replaced by a ninth fold held in reserve, so the eight-fold contrast survives. The
risk is worth taking because fewer folds cannot decorrelate the two candidate causes.

\startappendix
\begin{fundingtable}
No funding & & & & & \\
\end{fundingtable}

\end{document}
